\section{The PQID-as-AD Canonical Pattern}
\label{sec:ad-binding}

PQCrypt's existence as a sibling of \pqid{} and \pqhash{} in the same monorepo is not aesthetic --- it is structural. The canonical usage pattern is to bind a \pqid{} as the AEAD's associated data:

\begin{lstlisting}[language=Python]
record_id = pqid.new_hlc()                               # 41 bytes
envelope  = pqcrypt.encrypt(plaintext, recipient_pub,
                            associated_data=record_id)
# ... later, after delivery ...
plaintext = pqcrypt.decrypt(envelope, recipient_priv,
                            expected_ad=record_id)       # MUST match
\end{lstlisting}

\textbf{What this defends.} An attacker who reroutes the envelope under a different \pqid{} cannot decrypt it: AEAD tag verification fails (AD mismatch) and the decryption path returns the generic \texttt{DecryptionError} (\cref{sec:failure-modes}). The envelope is cryptographically bound to the identifier it was issued for.

\textbf{Cost.} Free. Every AEAD's AD path is computed at no marginal cost beyond the bytes being mixed into the MAC. PQCrypt's bench measurement \texttt{encrypt\_with\_pqid\_ad} confirms this experimentally (\cref{sec:appendices}).

\textbf{Cross-project AD patterns.} Other useful AD choices:
\begin{itemize}
\item A \pqhash{} of the plaintext, for envelopes that explicitly bind to content (defense against substituting one record's envelope onto another record).
\item A concatenation of a \pqid{} and a recipient identifier, for envelopes addressed to a specific (record, recipient) pair.
\item A protocol-level domain separator (e.g., \texttt{"pqcrypt-app-v1-records"}) plus the \pqid{}.
\end{itemize}

\textbf{Conformance.} The KAT file \texttt{pqcrypt/conformance/vectors/ad\_binding.json} exercises positive cases (envelope decrypts with the matching AD) and negative cases (any single-bit AD change causes a generic \texttt{DECRYPT\_FAILED}, with timing equal to the positive case to within the constant-time tolerance defined in \cref{sec:side-channels}).
